%% kapa: exam=1 round=34 year=2023 period=1 form=A subject=economics subject_ko=경제학원론 seq=21 number=61
%% source: Q-Net 감정평가사 기출문제 (cst00309) artlSeq=5211789 file=raw/1008/5211789/2023년도 제34회 감정평가사 1차시험 1교시.pdf
%% conversion: kordoc 4.24.51 -> tools/gen1.py (auto); review=proofread-sonnet; flags=none
\begin{question}{61}
고전학파의 국민소득결정모형에 관한 설명으로 옳지 않은 것은?
\begin{choices}
  \item 세이의 법칙(Say's law)이 성립하여, 수요측면은 국민소득 결정에 영향을 미치지 못한다.
  \item 물가와 임금 등 모든 가격이 완전히 신축적이고, 노동시장은 균형을 달성한다.
  \item 노동시장의 수요는 실질임금의 함수이다.
  \item 노동의 한계생산이 노동시장의 수요를 결정하는 중요한 요인이다.
  \item 통화공급이 증가하여 물가가 상승하면, 노동의 한계생산이 증가한다.
\end{choices}
\end{question}
\answer{5}
